\section{Preliminaries and Problem Formulation}

\subsection{Control Barrier Functions}

We consider a dynamical system $\dot{x} = f(x)$ where $x \in \R^n$ is the state and $f: \R^n \to \R^n$ is the drift dynamics. A set $\cS = \{x \in \R^n : h(x) \geq 0\}$ is forward invariant if there exists a Control Barrier Function (CBF) $h: \R^n \to \R$ satisfying the following condition for all $x \in \R^n$:

\begin{equation}
    \dot{h}(x) + \alpha(h(x)) \geq 0
    \label{eq:cbf_condition}
\end{equation}

where $\alpha: \R \to \R$ is a extended class-$\mathcal{K}$ function. For simplicity, we use the linear form $\alpha(h) = \alpha_0 h$ with $\alpha_0 > 0$. The time derivative of $h$ along the system dynamics is:

\begin{equation}
    \dot{h}(x) = \nabla h(x) \cdot f(x)
\end{equation}

where $\nabla h(x) \in \R^{1 \times n}$ is the gradient of $h$ with respect to $x$.

\subsection{CROWN Linearization for Neural Network CBFs}

When $h_\theta(x)$ is represented by a neural network with parameters $\theta$, directly verifying condition \eqref{eq:cbf_condition} is challenging due to the nonlinear activation functions. We employ CROWN (Complete ROugh WORth) linearization to compute rigorous lower and upper bounds on $h_\theta(x)$ and its derivatives over a region $\cR$.

For a simplicial region $\cS_\Delta$ with vertices $\{v_0, v_1, \ldots, v_n\}$ where each $v_i \in \R^n$, CROWN computes bounds of the form:

\begin{equation}
    A_L \leq \frac{\partial h_\theta}{\partial x} \leq A_U, \quad b_L \leq h_\theta(x) \leq b_U \quad \forall x \in \cS_\Delta
    \label{eq:crown_bounds}
\end{equation}

where $A_L, A_U \in \R^{n \times n}$ and $b_L, b_U \in \R^n$ are coefficient matrices and vectors that linearize the neural network.

\subsection{McCormick Relaxations for Product Bounds}

To bound the Lie derivative $\nabla h \cdot f$, we need to bound products of the form $A \cdot f$ where $A$ is the Jacobian bound and $f$ is the dynamics. We employ McCormick relaxations to compute:

\begin{equation}
    M_D \leq \text{diag}(A) \cdot f(x) \leq M_U \quad \forall x \in \cS_\Delta
\end{equation}

The McCormick lower bound for the product $y = w \cdot z$ with $w_L \leq w \leq w_U$ and $z_L \leq z \leq z_U$ is:

\begin{equation}
    y \geq w_L z + w z_L - w_L z_L + \eta((w_U - w_L)(z_U - z_L) - (w - w_L)(z - z_L))
    \label{eq:mccormick}
\end{equation}

where $\eta \in [0,1]$ is a relaxation parameter.

\section{Theoretical Framework}

\subsection{CBF Verification Over Simplicial Regions}

\begin{definition}[Simplicial Region]
    A simplicial region $\cS_\Delta \subset \R^n$ is defined by the convex hull of $n+1$ affinely independent vertices $\{v_0, v_1, \ldots, v_n\}$. Any point $x \in \cS_\Delta$ can be expressed as:
    \begin{equation}
        x = \sum_{i=0}^n \lambda_i v_i, \quad \sum_{i=0}^n \lambda_i = 1, \quad \lambda_i \geq 0
        \label{eq:barycentric}
    \end{equation}
    where $\lambda_i$ are the barycentric coordinates.
\end{definition}

\begin{assumption}
    The dynamics $f(x)$ is Lipschitz continuous on $\cS_\Delta$, and the neural network $h_\theta(x)$ is differentiable almost everywhere on $\cS_\Delta$ (satisfied by standard architectures with ReLU, Tanh, Sigmoid activations).
    \label{assum:lipschitz}
\end{assumption}

\begin{definition}[CBF Lower Bound on Simplicial Region]
    For a safe simplicial region $\cS_\Delta$, define the CBF lower bound:
    \begin{equation}
        \text{min\_L}(\cS_\Delta) = \min_{x \in \cS_\Delta} \left( \nabla h_\theta(x) \cdot f(x) + \alpha_0 h_\theta(x) \right)
        \label{eq:min_L_def}
    \end{equation}
    The region $\cS_\Delta$ is certified safe if and only if $\text{min\_L}(\cS_\Delta) \geq 0$.
\end{definition}

The minimum in \eqref{eq:min_L_def} can be computed efficiently since it is a linear function of $\lambda$ in barycentric coordinates. Specifically, substituting \eqref{eq:barycentric} into \eqref{eq:min_L_def} yields:

\begin{equation}
    \text{min\_L}(\cS_\Delta) = \min_{\lambda \geq 0, \mathbf{1}^T \lambda = 1} M^T \lambda + c
    \label{eq:min_L_bary}
\end{equation}

where $M$ and $c$ are coefficients from the CROWN linearization of $\nabla h_\theta \cdot f + \alpha_0 h_\theta$ over $\cS_\Delta$. This is a linear programming problem solvable in closed form.

\subsection{Randomized Smoothing Jacobian Estimator}

Computing exact gradients $\nabla_\theta h_\theta(x)$ through backpropagation is problematic when activation functions like ReLU have nondifferentiable points. We propose a Randomized Smoothing (RS) Jacobian estimator that avoids this issue.

\begin{definition}[Randomized Smoothing Jacobian]
    For a function $\psi: \R^p \to \R$ and parameter $\theta \in \R^p$, the RS Jacobian at $\theta$ with smoothing distribution $\epsilon \sim \Normal(0, \sigma^2 I_p)$ is:
    \begin{equation}
        J_{RS}(\theta) = \frac{1}{N \sigma^2} \sum_{i=1}^N \psi(\theta + \epsilon_i) \epsilon_i
        \label{eq:rs_jacobian}
    \end{equation}
    where $\epsilon_i \sim \Normal(0, \sigma^2 I_p)$ are i.i.d. samples.
\end{definition}

\begin{protocol}[RS Jacobian Estimation for CBF]
    \label{prot:rs_jacobian}
    For the CBF condition lower bound $\psi_s(\theta) = \text{min\_L}(\cS_\Delta; \theta)$ associated with simplicial $s$, the RS Jacobian over $S$ simplices is:
    \begin{equation}
        J_{RS}[s] = \frac{1}{N \sigma^2} \sum_{i=1}^N \psi_s(\theta + \epsilon_i) \epsilon_i
        \label{eq:rs_cbf_jacobian}
    \end{equation}
    where $N$ is the number of samples and $\sigma$ is the smoothing scale. The shared epsilon strategy computes this efficiently by perturbing $\theta$ once and evaluating all simplices in parallel.
\end{protocol}

\begin{property}[Unbiasedness of RS Jacobian]
    Under Assumption~\ref{assum:lipschitz}, the RS Jacobian estimator \eqref{eq:rs_jacobian} satisfies:
    \begin{equation}
        \E_\epsilon[J_{RS}(\theta)] = \nabla_\theta \E_\epsilon[\psi(\theta + \epsilon)] = \nabla_\theta \bar{\psi}(\theta)
        \label{eq:rs_unbiased}
    \end{equation}
    where $\bar{\psi}$ is the smoothed function. For small $\sigma$, $\bar{\psi}(\theta) \approx \psi(\theta)$.
\end{property}

\section{Constrained Gradient Optimization via QP Projection}

\subsection{The Repair Loss Function}

Given a set of failing simplicial regions, we define the repair loss to drive the network toward satisfying CBF conditions:

\begin{equation}
    \cL_{repair}(\theta) = \cL_{unsafe}(\theta) + \cL_{safe}(\theta)
    \label{eq:repair_loss}
\end{equation}

\begin{itemize}
    \item \textbf{Unsafe region violations} ($F_h$: $h(x) > 0$ in unsafe set):
    \begin{equation}
        \cL_{unsafe}(\theta) = \frac{1}{|F_h|} \sum_{x \in F_h} \text{softplus}(h_\theta(x) + \delta_h)
        \label{eq:loss_unsafe}
    \end{equation}
    \item \textbf{Safe region CBF violations} ($F_{safe}$: CBF condition violated):
    \begin{equation}
        \cL_{safe}(\theta) = \frac{1}{|F_{safe}|} \sum_{x \in F_{safe}} \text{softplus}(\delta_c - \nabla h_\theta(x) \cdot f(x) - \alpha_0 h_\theta(x))
        \label{eq:loss_safe}
    \end{equation}
\end{itemize}

where $\text{softplus}(z) = \frac{1}{\beta} \log(1 + e^{\beta z})$ with $\beta > 0$.

The gradient of \eqref{eq:repair_loss} is:
\begin{equation}
    g_F = \nabla_\theta \cL_{repair}(\theta)
    \label{eq:gf_grad}
\end{equation}

However, directly applying $g_F$ may violate the CBF conditions on already-verified safe regions. We must project $g_F$ onto a constraint set.

\subsection{QP Constraint Projection}

Let $J_{verified} \in \R^{M \times P}$ be the Jacobian matrix where row $j$ corresponds to the gradient of the CBF condition at the $j$-th verified simplicial region, and $P$ is the number of parameters. For safe regions, the constraint is:
\begin{equation}
    J_{safe}[j,:] \cdot d \geq 0 \quad \forall j \in \text{safe regions}
    \label{eq:constraint_safe}
\end{equation}
meaning that a parameter update direction $d$ should not decrease the CBF lower bound.

For unsafe regions where $h(x) > 0$, we require:
\begin{equation}
    J_{unsafe}[j,:] \cdot d \leq 0 \quad \forall j \in \text{unsafe regions}
    \label{eq:constraint_unsafe}
\end{equation}
meaning $h(x)$ should decrease.

Combining both, the constraint is:
\begin{equation}
    J_{verified} \cdot d \leq 0
    \label{eq:combined_constraint}
\end{equation}

We seek an update direction $d$ that is close to $-g_F$ (the gradient descent direction) while satisfying \eqref{eq:combined_constraint}. This yields the quadratic programming (QP) problem:

\begin{mini}
    {d \in \R^P}
    {\frac{1}{2} \|d + g_F\|^2}
    {\text{s.t.}}{J_{verified} \cdot d \leq 0}
    \label{eq:qp_primal}
\end{mini}

\begin{lemma}[QP Dual Problem]
    The dual of \eqref{eq:qp_primal} is:
    \begin{mini}
        {\lambda \in \R^M_{\geq 0}}
        {\frac{1}{2} \lambda^T (J_{verified} J_{verified}^T) \lambda - (J_{verified} g_F)^T \lambda}
        {}
        \label{eq:qp_dual}
    \end{mini}
    where $\lambda \in \R^M$ are the dual variables.
\end{lemma}

\begin{proof}
    The Lagrangian of \eqref{eq:qp_primal} is:
    \begin{equation}
        \cL(d, \lambda) = \frac{1}{2} \|d + g_F\|^2 + \lambda^T (J_{verified} d)
    \end{equation}
    with $\lambda \geq 0$. Setting $\nabla_d \cL = 0$:
    \begin{equation}
        d^* = -g_F - J_{verified}^T \lambda
    \end{equation}
    Substituting into $\cL$:
    \begin{align}
        \cL(\lambda) &= \frac{1}{2} \| -g_F - J_{verified}^T \lambda + g_F \|^2 + \lambda^T J_{verified}(-g_F - J_{verified}^T \lambda) \\
        &= \frac{1}{2} \lambda^T (J_{verified} J_{verified}^T) \lambda - \lambda^T J_{verified} g_F
    \end{align}
    which is exactly \eqref{eq:qp_dual}.
\end{proof}

\begin{algorithm}[t]
    \caption{QP Projection for Constrained Parameter Update}
    \label{alg:qp_projection}
    \begin{algorithmic}
        \REQUIRE Current parameters $\theta$, gradient $g_F$, verified Jacobian $J_{verified}$, learning rate $\eta$
        \STATE Compute normalized directions: $\hat{g} = g_F / \|g_F\|$, $\hat{J}[i,:] = J_{verified}[i,:] / \|J_{verified}[i,:]\|$
        \STATE Solve dual QP \eqref{eq:qp_dual} for $\lambda^*$ using projected gradient descent
        \STATE Compute update direction: $d = \hat{g} - \hat{J}^T \lambda^*$
        \STATE Update parameters: $\theta_{new} = \theta - \eta \cdot d$
        \RETURN Updated parameters $\theta_{new}$
    \end{algorithmic}
\end{algorithm}

\begin{property}[Constraint Satisfaction]
    The QP projection guarantees that for all verified simplices $j$:
    \begin{equation}
        J_{verified}[j,:] \cdot d^* \leq 0
        \label{eq:constraint_satisfied}
    \end{equation}
    meaning the update direction does not decrease any verified CBF lower bound.
\end{property}

\begin{property}[Descent Property]
    If $g_F \neq 0$ and the QP is feasible, the projected direction satisfies:
    \begin{equation}
        g_F \cdot d^* \leq 0
        \label{eq:descent}
    \end{equation}
    ensuring the loss non-increases along $d^*$.
\end{property}

\section{Progressive Depth Phased Repair Algorithm}

\subsection{Classification of Verification Failures}

During CBF verification with CROWN, a simplicial region $\cS_\Delta$ may fail verification in one of four ways:

\begin{enumerate}
    \item \textbf{$F_h$ (Unsafe-positive)}: $h(x) > 0$ in the unsafe set, i.e., $\max_{x \in \cS_\Delta \cap \text{unsafe}} h(x) > 0$
    \item \textbf{$F_{safe}$ (Safe CBF violation)}: CBF condition violated in safe region, i.e., $\text{min\_L}(\cS_\Delta) < 0$
    \item \textbf{$F_{depth}$ (Depth limit reached)}: Verification reached maximum depth without conclusive result
    \item \textbf{$F_{split}$ (Cannot split)}: Region too small to subdivide further, verification inconclusive
\end{enumerate}

The first two are \emph{definitive violations} (certainly unsafe), while the latter two are \emph{uncertain regions} that may or may not violate CBF conditions.

\begin{assumption}
    For a fixed max depth $D$, if a region is truly safe, then there exists some depth $d \leq D$ such that CROWN verification at depth $d$ will certify it.
    \label{assum:depth_soundness}
\end{assumption}

\subsection{Two-Phase Repair Strategy}

\begin{protocol}[Progressive Depth Phased Repair]
    \label{prot:phased_repair}
    The repair proceeds in two phases:

    \textbf{Phase 1 (Definitive Repair at Low Depth):}
    \begin{enumerate}
        \item Set $\text{max\_depth} = D_{start}$ (small, e.g., 10)
        \item Verify all regions and collect failures $\{F_h, F_{safe}, F_{depth}, F_{split}\}$
        \item Compute RS Jacobian $J_{RS}$ over verified regions
        \item Repair only definitive violations $\{F_h, F_{safe}\}$
        \item If $F_h = \emptyset$ and $F_{safe} = \emptyset$, transition to Phase 2
    \end{enumerate}

    \textbf{Phase 2 (Progressive Depth Increase):}
    \begin{enumerate}
        \item Increase $\text{max\_depth}$ to next value in schedule $\{D_{start}, D_2, D_3, \ldots\}$
        \item Re-verify all regions at new depth
        \item Repair all failures $\{F_h, F_{safe}, F_{depth}, F_{split}\}$
        \item Continue until $\text{max\_depth} = D_{limit}$ or all regions certified
    \end{enumerate}
\end{protocol}

\begin{property}[Soundness of Phased Repair]
    Under Assumption~\ref{assum:depth_soundness}, the progressive depth strategy never ``un-repairs'' a certified region: if a region is certified at depth $d$, it remains certified at any $d' \geq d$.
\end{property}

\begin{property}[Completeness with Depth]
    As $\text{max\_depth} \to \infty$, the verified region converges to the true safe set (up to splitting limit).
\end{property}

\begin{algorithm}[t]
    \caption{Progressive Depth Phased Repair}
    \label{alg:progressive_repair}
    \begin{algorithmic}
        \REQUIRE Initial model $h_\theta$, dynamics $f$, depth schedule $\cD = \{d_1, d_2, \ldots, d_K\}$
        \STATE $\theta \leftarrow$ initial parameters
        \STATE $phase \leftarrow 1$, $d \leftarrow d_1$
        \WHILE{not converged}
        \STATE $\{F_h, F_{safe}, F_{depth}, F_{split}\} \leftarrow$ verify$(\theta, f, d)$
        \IF{$phase = 1$ and $F_h = \emptyset$ and $F_{safe} = \emptyset$}
        \STATE $phase \leftarrow 2$, advance $d$ to next in $\cD$
        \ENDIF
        \IF{$phase = 1$}
        \STATE $\{F_h^{target}, F_{safe}^{target}\} \leftarrow \{F_h, F_{safe}\}$
        \ELSE
        \STATE $\{F_h^{target}, F_{safe}^{target}, F_{depth}^{target}, F_{split}^{target}\} \leftarrow \{F_h, F_{safe}, F_{depth}, F_{split}\}$
        \ENDIF
        \STATE $g_F \leftarrow$ compute\_repair\_gradient$(\theta, F_h^{target}, F_{safe}^{target})$
        \STATE $J_{RS} \leftarrow$ compute\_rs\_jacobian$(\theta, V_{safe}, V_{unsafe})$
        \STATE $d_{update} \leftarrow$ qp\_projection$(g_F, J_{RS})$
        \STATE $\theta \leftarrow \theta - \eta \cdot d_{update}$
        \ENDWHILE
        \RETURN Repaired model parameters $\theta$
    \end{algorithmic}
\end{algorithm}

\subsection{Stopping Criteria}

\begin{condition}[Convergence Criteria]
    The algorithm terminates when any of the following holds:
    \begin{enumerate}
        \item All regions certified: $|F_h| = |F_{safe}| = |F_{depth}| = |F_{split}| = 0$
        \item Maximum depth plateau: at $d = D_{limit}$, improvement $< \tau$ for $T$ consecutive iterations
        \item Global plateau: improvement $< \tau$ for $T$ consecutive iterations in Phase 2
        \item Maximum iterations reached
    \end{enumerate}
    where $\tau$ is the plateau threshold.
\end{condition}
